\documentclass[12pt,a4paper]{article}

% ─── Page Layout ───────────────────────────────────────────────────────────────
\usepackage[a4paper, margin=0.8in, headheight=14pt]{geometry}

% ─── Fonts (XeLaTeX) ───────────────────────────────────────────────────────────
\usepackage{fontspec}
\usepackage{unicode-math}
\setmainfont{TeX Gyre Pagella}
\setmonofont[Scale=0.88]{TeX Gyre Cursor}
\setmathfont{Latin Modern Math}

% ─── Packages ──────────────────────────────────────────────────────────────────
\usepackage{xcolor}
\usepackage{colortbl}
\usepackage{titlesec}
\usepackage{enumitem}
\usepackage{tabularx}
\usepackage{booktabs}
\usepackage{array}
\usepackage{amsmath}
\usepackage{tikz}
\usetikzlibrary{shapes.geometric, arrows.meta, positioning, calc, fit, decorations.pathreplacing, patterns}
\usepackage{tcolorbox}
\tcbuselibrary{skins, breakable}
\usepackage{hyperref}
\usepackage{fancyhdr}
\usepackage{graphicx}
\usepackage{microtype}
\hyphenpenalty=10000
\exhyphenpenalty=10000
\sloppy
\usepackage{caption}
\usepackage{setspace}
\usepackage{multirow}
\usepackage{tocloft}

% ─── Q&A helper ────────────────────────────────────────────────────────────────
\newcommand{\Q}[1]{\textbf{#1}\par\noindent\ignorespaces}

% ─── Colour Palette ────────────────────────────────────────────────────────────
\definecolor{myred}{RGB}{196,30,58}
\definecolor{mydark}{RGB}{30,30,50}
\definecolor{mygreen}{RGB}{14,120,80}
\definecolor{myteal}{RGB}{0,130,140}
\definecolor{myamber}{RGB}{180,100,0}
\definecolor{mypurple}{RGB}{100,30,160}
\definecolor{mygray}{RGB}{246,247,249}
\definecolor{mgframe}{RGB}{180,185,200}
\definecolor{watermark}{RGB}{145,150,170}
\definecolor{hdrred}{RGB}{250,232,235}
\definecolor{hdrgreen}{RGB}{220,244,233}
\definecolor{hdrteal}{RGB}{220,242,244}
\definecolor{hdramber}{RGB}{253,242,215}
\definecolor{hdrpurple}{RGB}{238,228,255}
\definecolor{hdrgray}{RGB}{238,240,245}
\definecolor{rowalt}{RGB}{250,251,253}

% ─── Section Formatting ────────────────────────────────────────────────────────
\titleformat{\section}{\large\bfseries\color{myred}}{\thesection.}{0.5em}{}[\vspace{1pt}{\color{myred}\rule{\linewidth}{0.8pt}}]
\titleformat{\subsection}{\normalsize\bfseries\color{myteal}}{\thesubsection}{0.5em}{}
\titlespacing*{\section}{0pt}{18pt}{7pt}
\titlespacing*{\subsection}{0pt}{11pt}{4pt}

% ─── Paragraph ─────────────────────────────────────────────────────────────────
\setlength{\parindent}{0pt}
\setlength{\parskip}{4pt}

% ─── TOC Styling ───────────────────────────────────────────────────────────────
\renewcommand{\cfttoctitlefont}{\large\bfseries\color{myred}}
\renewcommand{\cftaftertoctitle}{\par\noindent{\color{myred}\rule{\linewidth}{0.8pt}}}
\setlength{\cftbeforetoctitleskip}{0pt}
\setlength{\cftaftertoctitleskip}{6pt}
\renewcommand{\cftsecfont}{\normalsize\bfseries\color{mydark}}
\renewcommand{\cftsecpagefont}{\normalsize\bfseries\color{mydark}}
\renewcommand{\cftsecleader}{\cftdotfill{\cftdotsep}}
\setlength{\cftbeforesecskip}{7pt}
\renewcommand{\cftsubsecfont}{\small\color{mydark}}
\renewcommand{\cftsubsecpagefont}{\small\color{mydark}}
\setlength{\cftbeforesubsecskip}{3pt}
\setlength{\cftsubsecindent}{1.4em}

% ─── Table helpers ─────────────────────────────────────────────────────────────
\renewcommand{\arraystretch}{1.45}
\setlength{\tabcolsep}{6pt}
\arrayrulecolor{mydark}
\newcolumntype{B}[1]{>{\small\bfseries\raggedright\arraybackslash}m{#1}}
\newcolumntype{Y}{>{\small\raggedright\arraybackslash}X}
\renewcommand{\tabularxcolumn}[1]{m{#1}}

% ─── Header / Footer ───────────────────────────────────────────────────────────
\pagestyle{fancy}
\fancyhf{}
\fancyhead[L]{\small\color{myred}\textbf{OE-EC803A}\;\color{mydark}--- Internet of Things}
\fancyhead[R]{\small\color{myteal}\textbf{Module 4 Notes}}
\fancyfoot[C]{\small\color{mydark}\thepage}
\fancyfoot[R]{\small\color{watermark}\textit{\copyright\ Biraj Sarkar}}
\renewcommand{\headrulewidth}{0.5pt}
\renewcommand{\footrulewidth}{0.3pt}

% ─── Hyperref ──────────────────────────────────────────────────────────────────
\hypersetup{colorlinks=true, linkcolor=myteal, urlcolor=myteal,
  pdfauthor={Biraj Sarkar}, pdftitle={OE-EC803A --- Module 4 Notes}}

% ─── tcolorbox styles ──────────────────────────────────────────────────────────
\tcbset{
  tikzbox/.style={colback=mygray, colframe=mgframe, boxrule=0.6pt, arc=4pt,
    left=8pt, right=8pt, top=6pt, bottom=6pt,
    colbacktitle=mgframe, coltitle=mydark, fonttitle=\small\bfseries},
  defbox/.style={colback=hdrgreen, colframe=mygreen, boxrule=0.8pt, arc=4pt,
    left=9pt, right=9pt, top=6pt, bottom=6pt,
    colbacktitle=mygreen, coltitle=white, fonttitle=\small\bfseries},
  formulabox/.style={colback=hdrteal, colframe=myteal, boxrule=0.8pt, arc=4pt,
    left=9pt, right=9pt, top=7pt, bottom=7pt,
    colbacktitle=myteal, coltitle=white, fonttitle=\small\bfseries},
  examplebox/.style={colback=hdramber, colframe=myamber, boxrule=0.8pt, arc=4pt,
    left=9pt, right=9pt, top=6pt, bottom=6pt,
    colbacktitle=myamber, coltitle=white, fonttitle=\small\bfseries},
  masonbox/.style={colback=hdrpurple, colframe=mypurple, boxrule=0.8pt, arc=4pt,
    left=9pt, right=9pt, top=6pt, bottom=6pt,
    colbacktitle=mypurple, coltitle=white, fonttitle=\small\bfseries},
  exambox/.style={colback=hdrpurple, colframe=mypurple, boxrule=1pt, arc=5pt,
    left=10pt, right=10pt, top=8pt, bottom=8pt,
    colbacktitle=mypurple, coltitle=white, fonttitle=\bfseries,
    breakable, title after break={\textbf{Frequently Asked / Viva Questions (contd.)}}}
}
\tcbset{breakable}

% ─── TikZ styles ───────────────────────────────────────────────────────────────
\tikzset{
  block/.style={rectangle, draw=myteal, fill=hdrteal,
    text width=2.4cm, align=center, minimum height=0.9cm,
    font=\small, rounded corners=3pt, line width=0.7pt},
  arrow/.style={-Stealth, thick, color=mydark},
  line/.style={draw, thick, color=mydark}
}

\begin{document}

% ═══════════════════════════════════════════════════════════════════════════════
% TITLE BLOCK
% ═══════════════════════════════════════════════════════════════════════════════
\begin{center}
  \hfill \textit{Exam-Ready Short Notes} \\
  \vspace{6pt}
  {\LARGE\bfseries\color{myred} OE-EC803A --- Internet of Things}\\[5pt]
  {\normalsize\color{mydark} Semester VIII \quad$\bullet$\quad 3L:0T:0P \quad$\bullet$\quad 3 Credits}\\[3pt]
  {\normalsize\color{myteal}\textbf{Module 4: Prototyping Physical Design \& Enclosures}}\\[3pt]
  {\small\color{mydark} CGEC \;$|$\; B.Tech ECE (2023--27) \;$|$\; \textit{Biraj Sarkar}}
\end{center}
\vspace{2pt}
\noindent{\color{myred}\rule{\linewidth}{1.5pt}}
\vspace{2pt}

\hypersetup{linkcolor=mydark}
\tableofcontents
\hypersetup{linkcolor=myteal}
\newpage

% ===============================================================================
\section{The Physical Prototyping Process}
% ===============================================================================

Physical design defines the mechanical structure, look, and user feel of an IoT device. An enclosure must protect the internal electronics while providing access to pins, ports, and sensors.

\subsection{Prototyping Workflow Steps}
\begin{itemize}[leftmargin=*, itemsep=3pt]
  \item \textbf{Preparation:} Gathering circuit dimensions, mounting hole coordinates, and connectors.
  \item \textbf{Sketching \& Exploration:} Creating rapid 2D hand sketches to study button placements and accessibility.
  \item \textbf{Iteration:} Generating digital CAD models and manufacturing initial prototypes.
  \item \textbf{Testing:} Assembling the electronics inside the case to check for connector alignment, heat dissipation, and physical tolerances.
\end{itemize}

% ===============================================================================
\section{Prototyping Methods}
% ===============================================================================

Various low-fidelity and high-fidelity fabrication methods are used depending on the design stage.

\subsection{Non-Digital Prototyping (Low-Fidelity)}
Before using digital machines, designers build fast, manual models using cardboard, foam core, acrylic sheets, and modeling clay.
\begin{itemize}[leftmargin=*, itemsep=2.5pt]
  \item \textbf{Advantages:} Extremely low cost, zero machine wait time, and allows testing basic dimensions in a few minutes.
  \item \textbf{Drawback:} Low precision, fragile, and cannot be used for mechanical testing.
\end{itemize}

\subsection{Laser Cutting (2D Vector Fabrication)}
Laser cutting uses a high-power laser beam to cut or engrave sheet materials (acrylic, MDF, plywood, cardboard) based on 2D vector drawing profiles.
\begin{itemize}[leftmargin=*, itemsep=3.5pt]
  \item \textbf{Vector Layouts:} Drawn in programs like Inkscape or AutoCAD using closed paths.
  \item \textbf{Kerf:} The width of the material that is vaporized by the laser beam during cutting (typically $0.1\text{mm}$ to $0.3\text{mm}$). To create tight press-fit or finger joints, the design dimensions must compensate for this kerf allowance.
\end{itemize}

\subsection{3D Printing (Additive Manufacturing)}
\begin{tcolorbox}[defbox, title={3D Printing (Additive)}]
  \textbf{3D Printing} is an additive manufacturing process that builds 3D solid structures layer-by-layer from a digital CAD file. The dominant technology is \textbf{FDM (Fused Deposition Modeling)}, which extrudes molten thermoplastic filament (PLA, ABS) through a heated nozzle.
\end{tcolorbox}
Key 3D printing parameters that affect enclosure durability and print time:
\begin{itemize}[leftmargin=*, itemsep=3pt]
  \item \textbf{Layer Height:} Controls vertical resolution (typically $0.1\text{mm}$ to $0.3\text{mm}$). Smaller values give smoother surfaces but increase print time.
  \item \textbf{Infill Density:} The internal grid structure of the print, ranging from $0\%$ (hollow) to $100\%$ (solid). Enclosures typically use $15\%$ to $30\%$ infill.
  \item \textbf{Wall/Shell Thickness:} The thickness of the outer boundaries, defining enclosure strength and screw-mounting thread durability.
\end{itemize}

\subsection{CNC Milling (Subtractive Manufacturing)}
CNC (Computer Numerical Control) milling uses rotating cutting bits to carve away material from a solid block of plastic, wood, or metal (aluminum).
\begin{itemize}[leftmargin=*, itemsep=3pt]
  \item \textbf{Subtractive Concept:} The opposite of 3D printing. It starts with a block and subtracts material.
  \item \textbf{Applications:} Best for high-strength enclosures, metal heat-sink cases, and prototyping circuit boards (PCB routing).
\end{itemize}

\subsection{Repurposing and Recycling}
Hacking off-the-shelf plastic junction boxes, containers, or toys by drilling holes for cables and sensors.
\begin{itemize}[leftmargin=*, itemsep=2.5pt]
  \item \textbf{Advantages:} Provides immediate structural protection, is highly cost-effective, and reduces e-waste footprint during research stages.
\end{itemize}

\newpage

% ═══════════════════════════════════════════════════════════════════════════════
% QUICK REVISION PAGE
% ═══════════════════════════════════════════════════════════════════════════════
\begin{center}
  {\large\bfseries\color{myred} Quick Revision — Module 4}\\[3pt]
  {\small\color{mydark} Physical Design \& Prototyping $|$ OE-EC803A}
\end{center}
\noindent{\color{myred}\rule{\linewidth}{1.2pt}}
\vspace{8pt}

\begin{tcolorbox}[defbox, title={\textbf{One-Line Definitions}}]
\begin{itemize}[leftmargin=*, itemsep=3pt, topsep=0pt]
  \item \textbf{Kerf:} The thickness of the cut line or material slot vaporized by a laser cutter.
  \item \textbf{FDM:} A 3D printing process that builds objects by extruding layers of melted plastic filament.
  \item \textbf{Subtractive Manufacturing:} A fabrication process where material is carved away from a solid block.
  \item \textbf{Additive Manufacturing:} A fabrication process where material is joined layer-by-layer to build a structure.
  \item \textbf{Infill:} The internal supporting pattern printed inside a 3D object to save time and material.
\end{itemize}
\tcblower
Understanding the kerf is critical for press-fit joints where interlocking teeth must hold together without glue.
\end{tcolorbox}

\vspace{6pt}

\begin{tcolorbox}[exambox, title={\textbf{Frequently Asked / Viva Questions}}]
\begin{enumerate}[leftmargin=*, topsep=0pt, label=\textbf{Q\arabic*.}, itemsep=6pt]
  \item \Q{What is the difference between additive and subtractive manufacturing?}
        Additive manufacturing (e.g. 3D printing) builds parts layer-by-layer from nothing, producing zero waste and allowing complex internal cavities. Subtractive manufacturing (e.g. CNC milling) carves parts out of a solid block of raw material, producing waste chips but yielding high strength and support for metals.
        
  \item \Q{Explain the significance of ``kerf'' in laser cutting.}
        The kerf is the width of the cut slot created by the laser beam vaporizing material. If a laser has a kerf of $0.2\text{mm}$, cutting a $10.0\text{mm}$ square will result in a piece that is $9.8\text{mm}$ wide. Designers must compensate for this kerf to create tight, interlocking press-fit joints.
        
  \item \Q{Why is infill density in 3D printing rarely set to $100\%$?}
        A $100\%$ infill uses massive amounts of filament and takes a very long time to print. Structural strength gains level off significantly beyond $40\%$ infill. Enclosures are usually printed with $15\%$ to $30\%$ infill using internal patterns (like grid or gyroid) to maximize strength-to-weight ratio.
        
  \item \Q{What is Fused Deposition Modeling (FDM)?}
        FDM is the most common 3D printing technology. It melts a plastic filament (like PLA or ABS) and extrudes it through a small nozzle, moving in X, Y, and Z axes to lay down thin layers that fuse together as they cool.
        
  \item \Q{Why would a designer choose CNC milling over 3D printing for an industrial IoT gateway enclosure?}
        CNC milling allows using materials like structural aluminum. Aluminum enclosures provide superior mechanical protection, complete electromagnetic shielding (Faraday cage), and act as built-in heat sinks for high-power processors.
        
  \item \Q{What is the purpose of low-fidelity physical prototyping using cardboard?}
        Low-fidelity prototyping allows designers to physically hold and check the dimensions of a device layout within minutes, helping catch major interface errors (like port placement blockage or button hand-reach issues) before executing expensive machine runs.
        
  \item \Q{What is repurposing/recycling in physical design?}
        Repurposing refers to taking an existing mass-produced plastic box (like an electrical junction box or soap container) and drilling custom ports to house the IoT circuit. This saves cost, is weather-sealed, and reduces development time.
        
  \item \Q{What is SLA (Stereolithography) 3D printing, and how does it differ from FDM?}
        SLA uses a liquid photopolymer resin that is cured layer-by-layer using a UV laser beam. It yields extremely high resolutions and smooth surface finishes compared to FDM, but the materials are more fragile and require post-print chemical cleaning.
\end{enumerate}
\end{tcolorbox}

\vspace{6pt}

\begin{tcolorbox}[tikzbox, title={\textbf{Comparison Snapshot}}]
\textbf{3D Printing vs. CNC Milling vs. Laser Cutting:}\par\vspace{6pt}\noindent
{\renewcommand{\arraystretch}{1.35}%
\begin{tabularx}{\linewidth}{|B{3.2cm}|Y|Y|Y|}
\hline
\textbf{Feature} & \textbf{3D Printing (FDM)} & \textbf{CNC Milling} & \textbf{Laser Cutting} \\
\hline
Process & Additive (layer-by-layer). & Subtractive (carving). & 2D Cutting / Engraving. \\
\hline
Dimensions & Complex 3D geometries. & 3D shapes with tool limitations. & Flat 2D sheets (can fold/slot). \\
\hline
Materials & PLA, ABS, PETG. & Metals (Al, Steel), Wood, Acetal. & Acrylic, MDF, Plywood, Cardboard. \\
\hline
Wastage & Very Low (only support material). & High (residual chips and dust). & Low-Medium (sheet offcuts). \\
\hline
Setup Speed & Very Fast (load STL to printer). & Slow (requires CAM setups, tools). & Fast (load vector file). \\
\hline
\end{tabularx}}
\end{tcolorbox}

\end{document}
